\documentclass[11pt,a4paper]{article}

% ── Codificación, fuentes y lenguaje ───────────────────────────────────────────
\usepackage[utf8]{inputenc}
\usepackage[T1]{fontenc}
\usepackage[default,scale=0.95]{sourcesanspro}
\renewcommand{\familydefault}{\sfdefault}
\usepackage[spanish,es-noshorthands,es-tabla]{babel}

% ── Geometría y espaciado ─────────────────────────────────────────────────────
\usepackage[top=2.2cm, bottom=2.5cm, left=2.2cm, right=2.2cm, headheight=15pt]{geometry}
\usepackage{setspace}
\setstretch{1.18}
\usepackage{parskip}

% ── Matemáticas y unidades ────────────────────────────────────────────────────
\usepackage{amsmath}
\usepackage{amssymb}
\usepackage{siunitx}

% ── Circuitos ─────────────────────────────────────────────────────────────────
\usepackage[european, straightvoltages]{circuitikz}

% ── Tablas ────────────────────────────────────────────────────────────────────
\usepackage{booktabs}
\usepackage{tabularx}
\usepackage{array}
\usepackage{longtable}
\usepackage{colortbl}
\usepackage{enumitem}

% ── Colores, cajas e iconos ───────────────────────────────────────────────────
\usepackage[dvipsnames]{xcolor}
\usepackage{tcolorbox}
\tcbuselibrary{skins, breakable}
\usepackage{fontawesome5}
\usepackage{graphicx}

% ── Listados de código (dark-mode infográfico) ────────────────────────────────
%   Estilo base: fondo oscuro con números de línea. Los bloques lstlisting se
%   envuelven automáticamente en una caja tcolorbox redondeada.
\usepackage{listings}

% ── Secciones con barra de color (infografía) ────────────────────────────────
\usepackage{titlesec}
\titleformat{\section}
  {\normalfont\LARGE\bfseries\color{azulNoche}}
  {\colorbox{cyanNeon}{\color{fondoOscuro}\thesection}}{0.7em}{}
  [\vspace{2pt}\color{cyanNeon}\rule{\linewidth}{1.8pt}]
\titlespacing*{\section}{0pt}{24pt}{10pt}
\titleformat{\subsection}
  {\normalfont\large\bfseries\color{azulNoche}}
  {\textcolor{cyanNeon}{\thesubsection}}{0.7em}{}
  [\vspace{1pt}\color{grisLinea}\rule{\linewidth}{0.6pt}]
\titlespacing*{\subsection}{0pt}{14pt}{6pt}
\titleformat{\subsubsection}
  {\normalfont\normalsize\bfseries\color{azulNoche}}
  {\textcolor{cyanNeon}{\thesubsubsection}}{0.7em}{}
\titlespacing*{\subsubsection}{0pt}{10pt}{4pt}

% ── Cabeceras y pies de página ────────────────────────────────────────────────
\usepackage{fancyhdr}
\usepackage{lastpage}
\pagestyle{fancy}
\fancyhf{}
\renewcommand{\headrulewidth}{0pt}
\renewcommand{\footrulewidth}{0pt}
\fancyfoot[C]{%
  \begin{minipage}{\linewidth}
    \vspace{2pt}
    \color{cyanNeon}\rule{\linewidth}{1.2pt}\\[2pt]
    \centering\color{grisTexto}\footnotesize
    Página \thepage\ de \pageref{LastPage}
  \end{minipage}%
}

% ── Hiperenlaces ──────────────────────────────────────────────────────────────
\usepackage[colorlinks=true, linkcolor=azulNoche, urlcolor=cyanNeon]{hyperref}
\sloppy
\emergencystretch 3em

% ── Paleta de colores — tecnológico dark-mode ──────────────────────────────────
\definecolor{fondoOscuro}{HTML}{0D1B2A}%             Dark navy — fondo banda título
\definecolor{verdeTurquesa}{HTML}{004D40}%           Verde turquesa — bandas de marca
\definecolor{azulNoche}{HTML}{1B3A6B}%               Midnight blue — secciones, marcos
\definecolor{azulMedio}{HTML}{2A5298}%               Azul medio — variante de acento
\definecolor{cyanNeon}{HTML}{00C8E0}%                Cyan eléctrico — acentos primarios
\definecolor{verdeSignal}{HTML}{00B887}%             Verde señal — fortalezas, éxito
\definecolor{naranjaVivo}{HTML}{FF7043}%             Naranja vivo — mejoras, advertencias
\definecolor{rojoAlerta}{HTML}{E53935}%              Rojo alerta — errores
\definecolor{grisPapel}{HTML}{F4F6FA}%               Fondo tarjetas claras
\definecolor{grisLinea}{HTML}{C8D0E0}%               Bordes sutiles
\definecolor{grisTexto}{HTML}{6B7A99}%               Texto secundario
\definecolor{amarilloNota}{HTML}{FFB300}%            Notas / advertencias doradas
\definecolor{codigoFondo}{HTML}{1E2D3D}%             Fondo del bloque de código
\definecolor{codigoComentario}{HTML}{7A8FA6}%        Comentarios
\definecolor{codigoCadena}{HTML}{FF9D5C}%            Cadenas / strings
\definecolor{codigoKeyword}{HTML}{00C8E0}%           Palabras clave
\definecolor{codigoNumero}{HTML}{00E5A0}%            Números

% ── Banda de título: portada infográfica ──────────────────────────────────────
%   Diseño en 2 capas: banda de color (por defecto fondo oscuro) + franja de
%   acento cyan abajo. Uso: \bandaTitulo[color]{título}{subtítulo}.
%   El icono se puede cambiar con \renewcommand{\iconoBanda}{...} (FontAwesome).
\newcommand{\iconoBanda}{microchip}
\newcommand{\bandaTitulo}[3][fondoOscuro]{%
  \begin{tcolorbox}[
    enhanced,
    colback=#1, colframe=#1,
    boxrule=0pt, arc=8pt, outer arc=8pt,
    width=\linewidth,
    top=18pt, bottom=6pt, left=14pt, right=14pt,
    borderline south={3.5pt}{0pt}{cyanNeon},
  ]
    \begin{minipage}[c]{0.07\linewidth}
      \centering
      \textcolor{cyanNeon}{\fontsize{30}{30}\selectfont\faIcon{\iconoBanda}}
    \end{minipage}%
    \hspace{8pt}%
    \begin{minipage}[c]{0.88\linewidth}
      {\fontsize{20}{24}\selectfont\bfseries\color{white}#2}\par\vspace{5pt}%
      \textcolor{cyanNeon!80!white}{\small\faIcon{angle-right}~#3}%
    \end{minipage}
    \vspace{4pt}
  \end{tcolorbox}%
  \vspace{8pt}%
}

% ── Tarjetas de datos (infografía) ────────────────────────────────────────────
\tcbset{
  tarjetaDato/.style={
    enhanced, breakable,
    arc=6pt, outer arc=6pt,
    colback=grisPapel, colframe=azulNoche,
    boxrule=1pt,
    fonttitle=\bfseries\small, coltitle=white,
    attach boxed title to top left={yshift=-3mm, xshift=6mm},
    boxed title style={
      colback=azulNoche, colframe=azulNoche,
      arc=4pt, boxrule=0pt,
    },
    drop fuzzy shadow=azulNoche!30!white,
    top=8pt, bottom=8pt, left=10pt, right=10pt,
  },
  tarjetaIcono/.style={
    enhanced, breakable,
    arc=6pt, outer arc=6pt,
    colback=white, colframe=grisLinea, boxrule=0.8pt,
    fonttitle=\bfseries\small, coltitle=azulNoche,
    borderline west={3pt}{0pt}{cyanNeon},
    drop fuzzy shadow=azulNoche!25!white,
    top=6pt, bottom=6pt, left=10pt, right=10pt,
  },
  cajaTitulo/.style={
    enhanced, breakable,
    arc=6pt, outer arc=6pt,
    colback=azulNoche!10!grisPapel, colframe=azulNoche,
    boxrule=1pt,
    fonttitle=\bfseries, coltitle=white,
    attach boxed title to top left={yshift=-3mm, xshift=6mm},
    boxed title style={
      colback=azulNoche, colframe=azulNoche,
      arc=4pt, boxrule=0pt,
    },
    drop fuzzy shadow=azulNoche!25!white,
    top=8pt, bottom=8pt, left=10pt, right=10pt,
  },
  cajaContenido/.style={
    enhanced, breakable,
    arc=5pt, outer arc=5pt,
    colback=grisPapel, colframe=grisLinea,
    boxrule=0.8pt,
    fonttitle=\bfseries\small, coltitle=azulNoche,
    top=6pt, bottom=6pt, left=10pt, right=10pt,
  },
  cajaFortaleza/.style={
    enhanced, breakable,
    arc=6pt, outer arc=6pt,
    colback=verdeSignal!8!white, colframe=verdeSignal,
    boxrule=1pt,
    borderline west={4pt}{0pt}{verdeSignal},
    fonttitle=\bfseries\small, coltitle=verdeSignal!70!black,
    drop fuzzy shadow=verdeSignal!20!white,
    top=6pt, bottom=6pt, left=10pt, right=10pt,
  },
  cajaMejora/.style={
    enhanced, breakable,
    arc=6pt, outer arc=6pt,
    colback=naranjaVivo!8!white, colframe=naranjaVivo,
    boxrule=1pt,
    borderline west={4pt}{0pt}{naranjaVivo},
    fonttitle=\bfseries\small, coltitle=naranjaVivo!80!black,
    drop fuzzy shadow=naranjaVivo!20!white,
    top=6pt, bottom=6pt, left=10pt, right=10pt,
  },
  cajaRecomendacion/.style={
    enhanced, breakable,
    arc=6pt, outer arc=6pt,
    colback=cyanNeon!6!white, colframe=cyanNeon!60!azulNoche,
    boxrule=1pt,
    borderline west={4pt}{0pt}{cyanNeon},
    fonttitle=\bfseries\small, coltitle=azulNoche,
    drop fuzzy shadow=azulNoche!20!white,
    top=6pt, bottom=6pt, left=10pt, right=10pt,
  },
}

% ── Ícono + texto (encabezados de sección y elementos) ────────────────────────
\newcommand{\iconotexto}[2]{\textcolor{cyanNeon}{\faIcon{#1}}~#2}

% ── Listados de código: estilo dark + auto-envoltura ─────────────────────────
\lstdefinestyle{estiloCodigo}{
    backgroundcolor=\color{codigoFondo},
    basicstyle=\ttfamily\small\color{white},
    commentstyle=\color{codigoComentario}\itshape,
    stringstyle=\color{codigoCadena},
    keywordstyle=\color{codigoKeyword}\bfseries,
    numberstyle=\tiny\color{codigoComentario},
    numbers=left,
    numbersep=10pt,
    showstringspaces=false,
    breaklines=true,
    breakatwhitespace=false,
    tabsize=4,
    frame=none,
    captionpos=b,
    aboveskip=6pt,
    belowskip=4pt,
    xleftmargin=14pt,
    framexleftmargin=14pt,
    literate=
      {á}{{\'a}}1 {é}{{\'e}}1 {í}{{\'i}}1 {ó}{{\'o}}1 {ú}{{\'u}}1
      {Á}{{\'A}}1 {É}{{\'E}}1 {Í}{{\'I}}1 {Ó}{{\'O}}1 {Ú}{{\'U}}1
      {ñ}{{\~n}}1 {Ñ}{{\~N}}1 {ü}{{\"u}}1 {Ü}{{\"U}}1
      {¿}{{?`}}1 {¡}{{!`}}1
}
\lstdefinestyle{estiloPython}{
    style=estiloCodigo,
    language=Python,
}
\lstset{style=estiloCodigo}
\BeforeBeginEnvironment{lstlisting}{%
  \begin{tcolorbox}[
    enhanced, breakable,
    arc=6pt, outer arc=6pt,
    colback=codigoFondo, colframe=azulNoche,
    boxrule=1pt,
    drop fuzzy shadow=azulNoche!25!white,
    top=2pt, bottom=2pt, left=0pt, right=0pt,
  ]%
}
\AfterEndEnvironment{lstlisting}{\end{tcolorbox}}

% ── Cajas de contenido temáticas ──────────────────────────────────────────────
\newtcolorbox{cajaRecuerda}{
    enhanced, breakable,
    arc=6pt, outer arc=6pt,
    colback=azulNoche!10!grisPapel,
    colframe=azulNoche, boxrule=1pt,
    borderline west={4pt}{0pt}{cyanNeon},
    fonttitle=\bfseries\small\color{white},
    title={\faIcon{info-circle}~Recuerda},
    attach boxed title to top left={yshift=-3mm, xshift=6mm},
    boxed title style={colback=azulNoche, arc=4pt, boxrule=0pt},
    drop fuzzy shadow=azulNoche!25!white,
    left=10pt, right=10pt, top=8pt, bottom=8pt,
}

\newtcolorbox{cajaConcepto}{
    enhanced, breakable,
    arc=6pt, outer arc=6pt,
    colback=verdeSignal!8!white,
    colframe=verdeSignal, boxrule=1pt,
    borderline west={4pt}{0pt}{verdeSignal},
    fonttitle=\bfseries\small\color{white},
    title={\faIcon{lightbulb}~Concepto clave},
    attach boxed title to top left={yshift=-3mm, xshift=6mm},
    boxed title style={colback=verdeSignal!80!black, arc=4pt, boxrule=0pt},
    drop fuzzy shadow=verdeSignal!20!white,
    left=10pt, right=10pt, top=8pt, bottom=8pt,
}

\newtcolorbox{cajaEjemplo}{
    enhanced, breakable,
    arc=6pt, outer arc=6pt,
    colback=cyanNeon!5!white,
    colframe=cyanNeon!70!azulNoche, boxrule=1pt,
    borderline west={4pt}{0pt}{cyanNeon},
    fonttitle=\bfseries\small\color{fondoOscuro},
    title={\faIcon{code}~Ejemplo},
    attach boxed title to top left={yshift=-3mm, xshift=6mm},
    boxed title style={colback=cyanNeon, arc=4pt, boxrule=0pt},
    drop fuzzy shadow=azulNoche!20!white,
    left=10pt, right=10pt, top=8pt, bottom=8pt,
}

\newtcolorbox{cajaEjercicio}{
    enhanced, breakable,
    arc=6pt, outer arc=6pt,
    colback=azulMedio!7!white,
    colframe=azulMedio, boxrule=1pt,
    borderline west={4pt}{0pt}{azulMedio},
    fonttitle=\bfseries\small\color{white},
    title={\faIcon{pencil-alt}~Ejercicio},
    attach boxed title to top left={yshift=-3mm, xshift=6mm},
    boxed title style={colback=azulMedio, arc=4pt, boxrule=0pt},
    drop fuzzy shadow=azulNoche!20!white,
    left=10pt, right=10pt, top=8pt, bottom=8pt,
}

\newtcolorbox{cajaReto}{
    enhanced, breakable,
    arc=6pt, outer arc=6pt,
    colback=naranjaVivo!6!white,
    colframe=naranjaVivo, boxrule=1pt,
    borderline west={4pt}{0pt}{naranjaVivo},
    fonttitle=\bfseries\small\color{white},
    title={\faIcon{fire}~Reto},
    attach boxed title to top left={yshift=-3mm, xshift=6mm},
    boxed title style={colback=naranjaVivo, arc=4pt, boxrule=0pt},
    drop fuzzy shadow=naranjaVivo!20!white,
    left=10pt, right=10pt, top=8pt, bottom=8pt,
}

% ── Indicadores de dificultad ─────────────────────────────────────────────────
\newcommand{\facil}{\textcolor{verdeSignal}{\faIcon{circle}\,\textbf{Fácil}}}
\newcommand{\intermedio}{\textcolor{naranjaVivo}{\faIcon{circle}\,\textbf{Intermedio}}}
\newcommand{\dificil}{\textcolor{rojoAlerta}{\faIcon{circle}\,\textbf{Difícil}}}

% ─────────────────────────────────────────────────────────────────────────────

\begin{document}

\bandaTitulo[fondoOscuro]{IA en ingeniería con un esquema determinista de tres capas}{Por qué el 84 \% de un sistema con LLM no debería contener un LLM}

\vspace{4pt}
\begin{center}
\footnotesize\color{grisTexto}
\faIcon{user-circle}~ELECTRONICA\quad
\faIcon{calendar}~2026-10-05\quad
\faIcon{folder-open}~51 directivas\quad
\faIcon{cogs}~62 scripts de ejecución\quad
\faIcon{sitemap}~23 flujos\quad
\faIcon{book}~68 sesiones documentadas
\end{center}

\section{\iconotexto{layer-group}{1. La asimetría}}
Un modelo de lenguaje es un generador de texto plausible. La ingeniería no pide
plausibilidad: pide que el mismo input produzca el mismo output y que se pueda
auditar por qué. Esa asimetría no es una discusión filosófica, es medible, y se
mide en fallos reales con coste concreto.

\section{\iconotexto{layer-group}{2. Las tres capas}}

\begin{center}
\begin{tabular}{@{}lll@{}}
\toprule
\textbf{Capa} & \textbf{Artefacto} & \textbf{Responde} \\
\midrule
1 · Directiva & \texttt{directives/*.yaml} & ¿qué? \\
2 · Orquestación & \texttt{flujo\_*.py}, \texttt{mcp\_*.py} & ¿cuándo? \\
3 · Ejecución & \texttt{execution/*.py} & ¿cómo? \\
\bottomrule
\end{tabular}
\end{center}

\begin{cajaConcepto}
\textbf{La regla.} La lógica de negocio no se delega al modelo. Se le pide lo
único que hace bien --- producir texto --- y el resto se calcula en código.
\end{cajaConcepto}

\textbf{Capa 1, la directiva.} Es configuración, no conversación: se versiona,
se revisa con un diff y la escribe quien entiende el dominio.

\begin{lstlisting}[language=Python,caption={La frontera LLM, medida sobre el disco}]
# 62 scripts de ejecución
# 52 sin ninguna llamada a LLM   <- 84 por ciento
# 10 con LLM (fronteras declaradas, ficheros nombrados)
# el enrutador decide el tier localmente: 0 llamadas a la API, coste 0
\end{lstlisting}

\begin{cajaRecuerda}
\textbf{Criterio explícito, para que el número sea auditable.}
13 scripts referencian maquinaria LLM; se excluyen ['enrutador.py', 'generar\_charlada\_latex.py', 'llm\_client.py'] (cliente y enrutador: no invocan modelo).
\end{cajaRecuerda}

\section{\iconotexto{bug}{3. Tres fallos que costaron dinero}}

\subsection*{3.1 El modelo puntuaba; el programa suma}
El sistema devolvía \texttt{status: ok} y un JSON bien formado. Cada ítem
llevaba su denominador correcto. El total \emph{no era la suma de sus propias
partes}: el modelo había inventado el campo y la aritmética era suya.

La causa no era de formato sino de \textbf{ROL}: se le pidió al LLM el total.
Cambiar la prosa no lo arregla ---se probó---: los denominadores pasaron a ser
correctos mientras el total seguía sin cuadrar.

\begin{cajaEjemplo}
El arreglo fue dejarle \emph{una sola} tarea --- puntuar ítems --- y que el
programa sumara. El máximo pasó a ser un invariante: si la suma se pasa del tope,
se recorta \emph{y se avisa}. Cada parcial se limita a su propio peso. Y lo
ilegible se cuenta y se dice: un ítem sin parcial genera aviso, porque ignorarlo
en silencio hace creer que el alumno contestó lo que no contestó.
\end{cajaEjemplo}

\subsection*{3.2 La escala la declara el instrumento}
Tres preguntas de un punto cada una (escala 3) contestadas a medias. Sumar en
crudo daba \textbf{1.5/10} y nivel \emph{Insuficiente}: el alumno había
sacado el 50 \%.

La escala de un examen no es 10: es la que declara el examen. La nota
institucional sí es sobre 10, así que hace falta una conversión explícita.

Hubo además un \textbf{bug espejo}: al armar el denominador con lo que
devolvía el modelo, un criterio que el modelo se comía desaparecía del
denominador y todo lo demás se normalizaba \emph{hacia arriba}. El alumno
cobraba por un fallo del modelo. Con rúbrica, la escala la fija el archivo de
configuración y un criterio ausente vale 0 sobre esa escala.

\begin{cajaConcepto}
Ausencia de dato no es dato. Sin escala declarada no se publica nota: se publica
el motivo y la suma observada, que es información útil para el humano.
\end{cajaConcepto}

\subsection*{3.3 El modo JSON no es un esquema}
\texttt{response\_format=\{"type": "json\_object"\}} promete \textbf{una sola
cosa}: que la respuesta sea un objeto. No promete las claves, ni los tipos, ni
el vocabulario. El esquema del prompt es texto, y el texto es la capa más débil
del contrato.

El agravante fue más tonto: la instrucción \emph{nunca se envió}. El script
construía el mensaje de usuario y se olvidaba de adjuntar el esquema, así que el
modelo recibía literalmente \emph{devuelve SOLO el JSON} --- sin saber de qué
JSON --- y respondió con la primera forma que se le ocurrió. Y como
\texttt{json.loads} no tiene nada que decir sobre claves faltantes, la
inspección \emph{pareció} válida.

\begin{cajaEjemplo}
El arreglo no fue en el prompt: fue un \textbf{validador en el programa}, con
vocabulario cerrado para los campos enumerados y comprobación de tipos para el
resto. Y la clave que decide si el flujo es viable se volvió obligatoria:
confundir \emph{no saber} con \emph{no contestar} convierte una inspección
que no inspeccionó nada en un informe aparentemente confiable.
\end{cajaEjemplo}

\section{\iconotexto{balance-scale}{4. Una idea formal que se transfiere}}
La \textbf{álgebra de veredictos} compartida por las dos auditorías del
proyecto:\textbf{gana el peor estado, nunca un promedio}, lo no verificado
impide el verde, y se distinguen tres estados que la intuición colapsa en uno:
\emph{no instalado}, \emph{caído} y \emph{parado}.

El motivo por el que esto es una idea y no una función: un promedio con tres
estados en rojo y uno en amarillo se ve verde mientras algo crítico no se midió.
Y la segunda mitad de la regla es la que se olvida: una sola definición
compartida. Dos álgebras que divergen en silencio son la forma más cara de
equivocarse en un módulo cuya única razón de ser es la consistencia.

\section{\iconotexto{microchip}{5. Los límites}}
La capa 1 todavía no es verificable por la auditoría: es informativa. Un motor
de asistente rotativo no es reproducible. El hardware disponible no entrena
modelos. Un patrón cuya frontera se declara se enseña mejor que uno que se vende.

\section{\iconotexto{check-circle}{6. Seis pasos para adoptarlo}}
\begin{enumerate}
\item Escriba el SOP antes del código. Si no sabe escribirlo, aún no sabe qué
      tiene que automatizar.
\item Ponga el criterio en el programa: toda aritmética, comparación o
      conversión vive en código, no en el prompt.
\item Valide la forma, no confíe en ella.
\item Defina la escala en el instrumento, y si no la hay no publique el número.
\item Verifique por registro, no por bandera verde: un programa que termina con
      código 0 puede estar mintiendo, así que hay que leer su salida.
\item Deje la frontera escrita: los ficheros con LLM se nombran.
\end{enumerate}

\vfill
\begin{center}
\footnotesize\color{grisTexto}
Este material lo genera el propio patrón del que habla: directiva, orquestador y
script, con la paleta de color compartida del workspace.
\end{center}

\end{document}
